% =============================================================================
% §1 Introduction (~1.0 page) — rewritten as progressive narrative
% =============================================================================
\section{Introduction}
\label{sec:intro}

Large language models increasingly act as autonomous agents that
orchestrate long-horizon executable tasks rather than answer
isolated queries \citep{Yao2023ReAct,Shinn2023Reflexion,Park2023Generative}.
Yet on composite real-world workloads, even the strongest
frontier systems remain measurably bounded
\citep{Merrill2026TerminalBench,OpenAIGDPVal2025}.  This
gap is difficult to close by raw model scaling alone.  A growing
consensus instead treats agent capability as a layered system, in
which a frozen foundation model is paired with externalised
skills and an agent harness that loads and executes them
\citep{AnthropicSkillCreator2025}.  A skill
packages reusable procedural knowledge invoked at inference time.
Unlike opaque, slow-to-update model weights or transient
in-context examples, it is interpretable and modular.  
The same externalise-and-retrieve pattern is already well-established for tool use,
where agents retrieve and invoke calls over large API inventories
\citep{Qin2024ToolLLM,Patil2023Gorilla} and the recent Model Context Protocol (MCP) further standardises access to external tools and resources \citep{Hou2025MCP,MCPCorpus2025}.
Skills extend this pattern from external tool calls to
procedural knowledge itself.

Industry-grade harnesses now ingest skill files routinely
\citep{AnthropicClaudeCode2025,OpenAIClaudeCode2025}, and a
community ecosystem has formed around the \texttt{SKILL.md}
specification \citep{AnthropicSkillCreator2025}, now spanning
thousands of public repositories.  Whether this abundance actually improves agent performance on real tasks, and what it takes to make it do so, remains
unsettled.  Curated skills paired one-to-one with tasks raise pass
rates on average but with sharply heterogeneous per-domain effects
\citep{SkillsBench2026}, while uncurated community libraries can
fail to improve over a no-skill baseline
\citep{SkillFlow2026,UCSBSkillsWild2026}.  Two gaps stand in the
way.  \textbf{First}, we know of no released corpus that is at once broadly
aggregated, strictly quality-gated, and cleared for open
redistribution.  Existing
artifacts draw from only one or two source channels~\citep{SkillFlow2026,UCSBSkillsWild2026}, and those that filter for quality show no evidence of broad coverage at
scale~\citep{SkillNet2026}.  
\textbf{Second}, the downstream value of a community corpus on real agent tasks is largely uncharacterised. Prior evaluations pair tasks with hand-picked oracle skills rather than a deployed corpus~\citep{SkillsBench2026}, run in simulated environments~\citep{SkillNet2026}, or measure in-the-wild usage through a single harness over limited sources~\citep{UCSBSkillsWild2026}.

% Closing both gaps calls for four capabilities at once.  Broad
% aggregation and strict curation address the resource side.
% Precise matching and end-to-end evaluation address the
% measurement side, since a fair test must retrieve skills from the
% corpus rather than hand-pick an oracle.  We present
% \textbf{SkillCorpus}, which delivers all four and closes both gaps.
To close both gaps, we present \textbf{SkillCorpus}, a single end-to-end framework that consolidates the open \texttt{SKILL.md} ecosystem into a curated, deployable corpus and evaluates it on real-world agent tasks — making it possible to measure when community skills help real agents, and where they do not.

Our contributions are threefold:
\begin{itemize}
\item \textbf{Framework.} We unify aggregation, curation, matching, and real-world evaluation of the open \texttt{SKILL.md} ecosystem, turning a fragmented, redundant, and uneven pool of community skills into a curated foundation for reuse at scale.

\item \textbf{Open resource.} We release the resulting $96{,}401$-skill corpus (curated from $\sim$$821{,}000$ raw files), its fine-tuned retrieval-and-selection stack, and the full curation, training, and evaluation code, all licence-audited and $100\%$ OSI-permissive — a fully open, ecosystem-scale resource for skill-augmented agent research.

\item \textbf{Findings.} Evaluating across three real-world agent benchmarks, independent harnesses, and open-to-frontier backbones (\S\ref{sec:experiments}), we find that integrating SkillCorpus consistently improves agent performance, with gains modulated by corpus coverage and the harness, characterising the real utility and limits of community skills.

%\item We propose SkillCorpus, a single end-to-end framework that
%      unifies all four capabilities over the open \texttt{SKILL.md}
%      ecosystem, turning a scattered, uneven pool of community skills
%      into a curated foundation for reuse at scale.

% \item We release the resulting $96{,}401$-skill corpus curated from
%       $\sim$$821{,}000$ raw files, its fine-tuned
%       retrieval-and-selection stack, and the full curation, training,
%       and evaluation code, all $100\%$ OSI-permissive,
%       an end-to-end open, ecosystem-scale resource for skill-augmented
%       agent research.

% \item Evaluating across multiple real-world agent benchmarks,
%       independent harnesses, and open-to-frontier backbones
%       (\S\ref{sec:experiments}), we find that integrating SkillCorpus
%       consistently improves agent performance, with gains modulated by
%       corpus coverage and the harness, characterising the real utility
%       and limits of community skills.
\end{itemize}
